\section{Discussion}
\label{sec:discussion}

The relational model for \ematching is
 not only simple and fast,
 but it opens the door to a wide range of future work
 to further improve \egraphs and \ematching.

\subsection{Relational Representation of Code}

\update{
The representation of \enodes as relations and patterns as
 relational queries is simple and natural.
It may even feel obvious to someone familiar with Prolog-style programming,
 where a function with arity $k$ is often written as a relation with arity $k+1$, 
 with 1 extra argument for the output.
This relational representation can also be found in the literature on congruence closure.
For example, \citet{rummer2012matching} uses the same encoding to 
 simulate congruence closure procedures with a hyper-resolution calculus.
Unlike our focus on the performance of \ematching, 
 their goal is to improve the completeness of quantified reasoning in SMT solvers.
Research on large-scale code search and analysis~\cite{urma2013expressive, avgustinov2016ql, doop, Glean}
 has also explored storing program repositories in a database.
Programmers may issue queries against the database 
 to find code that match certain patterns, 
 which may be examples of API usage or 
 ``anti-patterns'' that pose security risks.
Certain complex queries can even express sophisticated analyses like 
 the points-to analysis~\cite{doop, avgustinov2016ql}. 
These code search and analysis engines need to balance the 
 need for expressiveness, efficiency, and ease of use.

While the idea to represent code as relations is not new, 
 we are the first, to the best of our knowledge, 
 to leverage relational join algorithms to speed up \ematching.
We designed optimizations specialized for \ematching, 
 and also derived the first non-trivial complexity bound
 for \ematching from a careful analysis of the generic join algorithm.
}
\subsection{Pushdown Optimization}

An \ematching pattern may have additional filtering
 condition associated with it.
For example, a rewrite rule with left hand side \verb|(* (/ x y) y)| may
 additionally require that $y\neq 0$.
When the variables involved in conditions all occur in a single relation (e.g.,
 $R_*$ and $R_/$),
 this relation can be effectively filtered
 even before being joined
 (e.g., using predicate $\find(\sigma(y))\neq\find(\lookup(0))$),
 which could immediately prune a large search space.

We call this \textit{pushdown optimization},
 which can be considered as
 \egraph's version of the relational query
 optimization that always pushes the filter operations down to the bottom
 of the join tree.
Note
 that the ability to do pushdown optimization stems
 from relational \ematching's ability to consider the constraints in any order;
 backtracking \ematching could not support this technique.
We currently do
 not implement this,
 because it requires breaking
 changes to \egg's interface.

Conditions that involve multiple variables can be ``pushed down'' as well.
In generic join,
 the filter can occur immediately after the variables appear in the variable ordering.
Thus,
 an implementation that supports these conditional filters
 should take this into account when determining variable ordering.

% The ``top-down'' approach of backtracking \ematching
%  cannot pre-filter
% To see this, consider the pattern $f(g(\alpha),g(\beta))$,
%  with the requirement $\sigma(\alpha)\neq 0$ and $\sigma(\beta)=0$.
% The \egraph cannot prune \enodes $g(i)$ using both
%  filtering condition $i=\find(\lookup(0))$
%  and $i\neq \find(\lookup(0))$ simultaneously.

\subsection{Join Algorithms}
Research in databases has proposed a myriad of different join algorithms.
For example, state-of-the-art database systems implement two-way joins
 like hash join and merge-sort join.
They have a longer history than generic join,
 and benefit from various constant factor optimizations.
Extensive research has focused on generating highly efficient query plans
 using two-way joins.
On the other hand, Yannakakis' algorithm~\cite{yannakakis}
 is proven to be optimal on a class of queries called {\em full ayclic queries},
 running in time linear to the total size of the input and the output.
All linear patterns correspond to acyclic queries,
 but some nonlinear patterns correspond to cyclic ones.
 Recent research~\cite{mhedhbi19,
 freitag} has also experimented with combining
 traditional join algorithms with generic join,
 achieving good performance.
In this paper we choose generic join for its simplicity,
 and future work may consider other join algorithms for relational \ematching.

\subsection{Incremental Processing}\label{sec:incremental}
We have focused on improving the core \ematching\ algorithm in this paper,
 yet prior work has successfully sped up \ematching\ by making it incremental~\cite{efficient-ematching}.
When the changes to the \egraph\ are small and the results of \ematching\
 patterns are frequently queried,
 maintaining the already-discovered matches becomes crucial for efficiency.
From our relational perspective, incremental \ematching\
 is captured precisely by the classic problem of
 incremental view maintenance (IVM)~\cite{DBLP:conf/vldb/CeriW91,
 DBLP:conf/sigmod/SalemBCL00,DBLP:conf/sigmod/ZhugeGHW95} in databases.
IVM aims to efficiently update an already-computed query result upon changes
 to the database, without recomputing the query from scratch.
\update{
Future research can follow our approach but implement e-graphs directly on top of
 a relational database engine, 
 inheriting the incrementality from the underlying system. 
For example, a datalog engine provides semi-naive evaluation.}

There is opportunity to improve relational \ematching\ even without a
 full-fledged IVM solution.
For example, we have shown in \autoref{fig:speedup} that index building
 can take up a significant portion of the run time.
Our index implementation is based on a hash trie,
 which is simple but difficult to update efficiently.
We are experimenting with an alternative index design based on sort tries,
 in the hope that it can make updates as simple as inserting into a sorted array.

\subsection{Building on Existing Database Systems}
Given our reduction from \ematching\ to conjunctive query answering,
 one may wonder if other \egraph operations could be reduced to
 relational operations so that a fully functioning \egraph engine
 can be implemented purely on top of an off-the-shelf database system.
There are many benefits to it.
For example, we could enjoy an industrial-strength query optimization
 and execution engine for free
 (although most industrial databases do not use worst-case optimal join algorithms),
 and eliminates the cost of transforming an \egraph
 to a relational database.
Moreover, this approach would enjoy any properties of the host database
 system, including persistence, incremental maintenance, concurrency, and fault-tolerance.
% {\color{blue} (It may also be worth mentioning the idea of scalability / persistent e-graphs / egraph as a
%  service / etc.)}

As a proof of concept, we implemented a prototype \egraph implementation on top of SQLite,
 an embedded relational database system, with 160 lines of Racket code.
\Egraph operations like insertion and merging are translated into
 high-level SQL queries and executed
 using SQLite.
% Therefore, the costs of rebuilding both the relational database and its indices
%  are naturally avoided,
%  which can sometimes be the bottleneck.
This na\"ive prototype is not competitive with highly optimized
 implementations like \egg,
 especially given our relational \ematching approach.
However, with appropriate indices and query plan,
 it could have similar asymptotic performance.
 % the standard \egraph algorithm.
Specialized data structures to represent equivalence relations~\cite{nappa2019fast}
 could also help performance.
Therefore, not only \ematching but also other \egraph operations can be
 expressed as relational queries,
 which hints at the possibility of developing real-world \egraph engines
 on top of existing relational database systems.

% \subsection{Other Potential Optimizations}

% \yz{I'm stuck on how to write this subsection.}

% There are many more techniques in the database community that we could use to
% achieve better performance but we have not implemented. In this subsection, we
% name a few potential optimizations that we are aware of that could potentially
% improve the performance of relational \ematching.

% \subsubsection*{Query Optimizations}

% \yz{I don't know much about query optimization literatures}

% There are many existing literatures that
% explores various query optimization techniques. For example, cardinality
% estimation is widely used in commercial databases to find efficient query plans.
% Such techniques could also be used in the setting of relational \ematching.

% \yz{
%   The emptyheaded paper also mentions hypertree decomposition as an optimization
%   technique for generic join, which is something we can mention.
% }

% \subsubsection*{Set Intersections}

% A key operation of generic join is set intersection. Currently we implement set
% intersection as simple hash lookup, which is nonetheless cache unfriendly.
% However, many literature has explored faster and more cache-friendly algorithms
% for set intersection. TODO.

% \yz{
%   There is a paper on set intersection (not sure if it's sota; also see related works): https://arxiv.org/pdf/1103.2409.pdf.
%   EmptyHeaded also describes several set intersection algorithms they use for generic
%   join that is optimized for SIMD.
% }

% \subsubsection*{In-memory Databases}

% \yz{we should probably mention that we
%   specialize our operators, which is great, but there are also tons of
%   researches done on speeding up in-memory databases that we can exploit.}

% \subsubsection*{Functional Dependencies}

% \subsection{Symmetry Breaking}

%%% Local Variables:
%%% mode: latex
%%% TeX-master: "main"
%%% End:
